\documentclass[aspectratio=169]{beamer}
\usepackage[utf8]{inputenc}
\usepackage{amsmath, amssymb}
\usepackage{graphicx}
\usepackage{booktabs}
\usepackage{tcolorbox}

% Theme styling based on UNIB colors
\definecolor{UNIBBlue}{RGB}{0, 51, 102}
\definecolor{UNIBYellow}{RGB}{255, 204, 0}
\setbeamercolor{palette primary}{bg=UNIBBlue,fg=white}
\setbeamercolor{palette secondary}{bg=UNIBYellow,fg=black}
\setbeamercolor{structure}{fg=UNIBBlue}
\setbeamercolor{title}{fg=UNIBBlue}
\setbeamercolor{frametitle}{bg=UNIBBlue,fg=white}

\title{Optimisasi untuk Teknik Elektro}
\subtitle{Minggu 3: Metode Simpleks \& Dualitas}
\author{Ir. Novalio Daratha S.T., M.Sc., Ph.D.}
\date{Semester Ganjil 2026}

\begin{document}

\begin{frame}
    \titlepage
\end{frame}

\begin{frame}{Tujuan Pembelajaran (CPMK-1, CPMK-3)}
    \begin{itemize}
        \item Mampu merubah bentuk standar program linier dengan variabel slack/surplus.
        \item Mampu menjelaskan konsep dan iterasi Metode Simpleks secara konseptual.
        \item Mampu mendefinisikan masalah dual dari masalah primal program linier.
        \item Mampu menginterpretasikan nilai dual secara ekonomi dalam konteks keteknikan.
    \end{itemize}
\end{frame}

\begin{frame}{Bentuk Standar Program Linier}
    \textbf{Semua kendala harus berupa persamaan ($=$) dengan ruas kanan non-negatif.}
    \vspace{0.5cm}
    \begin{itemize}
        \item Kendala $\le$ (kurang dari atau sama dengan) ditambahkan \textbf{Variabel Slack} (waktu/sumber daya yang tersisa).
        \item Kendala $\ge$ (lebih dari atau sama dengan) dikurangi \textbf{Variabel Surplus} (kelebihan di atas minimum).
    \end{itemize}
\end{frame}

\begin{frame}{Contoh Bentuk Standar}
    \textbf{Primal:}
    Maksimumkan $Z = 3x_1 + 5x_2$ \\
    s.t. \\
    $x_1 \le 4$ \\
    $2x_2 \le 12$ \\
    $3x_1 + 2x_2 \le 18$ \\
    $x_1, x_2 \ge 0$
    
    \vspace{0.5cm}
    \textbf{Bentuk Standar:}
    Maksimumkan $Z = 3x_1 + 5x_2 + 0s_1 + 0s_2 + 0s_3$ \\
    s.t. \\
    $x_1 + s_1 = 4$ \\
    $2x_2 + s_2 = 12$ \\
    $3x_1 + 2x_2 + s_3 = 18$ \\
    $x_1, x_2, s_1, s_2, s_3 \ge 0$
\end{frame}

\begin{frame}{Metode Simpleks (Konseptual)}
    \begin{itemize}
        \item Pendekatan aljabar untuk mencari solusi optimal pada titik-titik sudut (extreme points) ruang solusi yang feasible.
        \item Iterasi bergerak dari satu solusi dasar feasible (BFS) ke BFS tetangganya yang memberikan nilai fungsi tujuan yang lebih baik.
        \item \textbf{Langkah Utama:}
        \begin{enumerate}
            \item Tentukan BFS awal (biasanya di titik asal).
            \item Cek optimalitas (Apakah ada variabel non-basis yang jika dimasukkan ke basis akan memperbaiki Z?).
            \item Tentukan variabel masuk (entering variable) dan variabel keluar (leaving variable).
            \item Lakukan operasi baris dasar (pivoting) untuk mendapatkan matriks BFS baru.
            \item Ulangi hingga kondisi optimal tercapai.
        \end{enumerate}
    \end{itemize}
\end{frame}

\begin{frame}{Dualitas (Duality)}
    \textbf{Setiap masalah program linier (Primal) memiliki masalah terkait yang disebut Dual.}
    \vspace{0.3cm}
    \begin{itemize}
        \item Solusi optimal Primal dan Dual memiliki nilai fungsi tujuan yang sama (Strong Duality).
        \item Memberikan perspektif ekonomi pada masalah.
        \item Variabel Dual ($y_i$) merepresentasikan \textit{shadow price} atau nilai batas dari sumber daya (kendala $i$).
    \end{itemize}
\end{frame}

\begin{frame}{Interpretasi Ekonomi Dualitas di Teknik Elektro}
    Contoh: Optimisasi alokasi daya dari beberapa generator untuk memenuhi beban (Primal: Minimisasi Biaya).
    \vspace{0.5cm}
    \begin{tcolorbox}[colback=UNIBYellow!20,colframe=UNIBBlue,title=Makna Variabel Dual (Shadow Price)]
        Berapa banyak \textbf{penghematan biaya} (atau pengurangan total objektif) yang bisa didapatkan jika batas kendala kapasitas daya diperlonggar sebesar 1 unit (1 MW)?
    \end{tcolorbox}
    \begin{itemize}
        \item Jika generator A sudah pada kapasitas maksimum, shadow price-nya positif (ada manfaat ekonomi jika kapasitas ditambah).
        \item Jika generator B belum mencapai batas, shadow price-nya nol (menambah batas kapasitas B tidak mengubah solusi optimal).
    \end{itemize}
\end{frame}

\begin{frame}{Kesimpulan}
    \begin{itemize}
        \item Bentuk standar (slack/surplus) diperlukan untuk algoritma aljabar seperti Simpleks.
        \item Metode Simpleks mencari solusi efisien di titik-titik sudut.
        \item Dualitas memberikan nilai "harga bayangan" (shadow price) untuk setiap sumber daya atau batas operasional, sangat berharga dalam keputusan investasi keteknikan.
    \end{itemize}
\end{frame}

\end{document}
